\documentclass[11pt]{article}
\usepackage[margin=2cm]{geometry}
\usepackage[utf8]{inputenc}
\usepackage[T1]{fontenc}
\usepackage{mathptmx}
\usepackage{enumitem}
\usepackage{booktabs}
\usepackage{xcolor}
\usepackage{hyperref}
\usepackage{amssymb}
\usepackage{setspace}
\onehalfspacing

\hypersetup{colorlinks=true,linkcolor=blue!70!black,urlcolor=blue!70!black}

\newcommand{\yes}{\rlap{$\checkmark$}\square}
\newcommand{\no}{\square}

\pagestyle{plain}

\begin{document}

\begin{flushright}
{\large\bfseries\color{black!70} nature\\portfolio}
\end{flushright}

\noindent Corresponding author(s): Lucas Rover\\
Last updated by author(s): 11-03-2026

\begin{center}
{\LARGE Machine Learning Checklist \normalsize V\,1.1}
\end{center}

\noindent Nature Portfolio wishes to improve the reproducibility of the work that we publish. This form is intended to provide structure for consistency and transparency in reporting of works using or developing Machine Learning models. Some list items might not apply to an individual manuscript, but all fields must be completed for clarity.

\bigskip
\hrule
\bigskip

%%=============================================================
\section*{1. Availability and reproducibility of Code and Data}

\noindent Please select all that apply regarding the availability of the data and code used in the study.

\begin{itemize}[leftmargin=2em]
\item[$\no$] Code will be included in a CodeOcean capsule.

\item[$\yes$] The \textbf{source code} is included in the submission or available in a public repository:\\
\textbf{URL:} \url{https://github.com/Roverlucas/genai-reproducibility-protocol}\\
{\small (Available to reviewers during peer review; publicly released upon publication.)}

\item[$\no$] A \textbf{compiled standalone version} of the software is included in the submission or available in a public repository.\\
{\small N/A --- Python scripts, no compilation needed.}

\item[$\yes$] A \textbf{test dataset} and instructions/scripts for replicating the results are included in the submission or available in a public repository:\\
\textbf{URL:} \url{https://github.com/Roverlucas/genai-reproducibility-protocol}\\
{\small \texttt{data/inputs/} contains all 30 abstracts and RAG contexts; run scripts reproduce all 4,104 experiments.}

\item[$\yes$] A \textbf{Readme file} with instructions for installing and running the code is included in the submission or available in a public repository:\\
\textbf{URL:} \url{https://github.com/Roverlucas/genai-reproducibility-protocol}\\
{\small README.md in repository root with full installation and execution instructions.}

\item[$\yes$] The code is made available to reviewers during review.

\item[$\yes$] \textbf{Pretrained models} are used in the study and accessible through:\\
\textbf{URL:} Ollama Hub (\url{https://ollama.com/library}): llama3:8b, mistral:7b, gemma2:9b\\
{\small Open-weight models, freely downloadable. API models (GPT-4, Claude, Gemini, DeepSeek, Perplexity) accessible via respective commercial APIs.}

\item[$\no$] \textbf{Pretrained models} are used in the study and are not accessible.

\item[$\yes$] The paper contains information on how to obtain code and data after publication.\\
{\small Data Availability and Code Availability statements specify Zenodo archival with persistent DOI upon publication.}
\end{itemize}

%%=============================================================
\section*{2. Datasets}

\begin{enumerate}[label=\Alph*.,leftmargin=2em]
\item All data sources are listed in the paper.\\
$\yes$ \textbf{Yes} \quad $\no$ No\\
{\small Methods \S Input data; Supplementary Section S2 (Table S1 lists all 30 abstracts); Supplementary Section S3 (RAG contexts).}

\item The train, test and validation datasets are publicly available, and links/accession numbers have been provided in the manuscript or supplementary materials.\\
$\yes$ \textbf{Yes} \quad $\no$ No\\
{\small All 30 input abstracts available at \texttt{data/inputs/abstracts.json} in the repository. No train/test/validation split --- this is an inference reproducibility study, not a training study.}

\item We have reported and discussed potential dataset biases in the paper. Where applicable, appropriate mitigation strategies were used.\\
$\yes$ \textbf{Yes} \quad $\no$ No\\
{\small Discussion \S Limitations: corpus limited to 30 English-language AI/ML abstracts; other domains and languages may show amplified effects. Balanced 10-abstract subsample analysis (Extended Data Table 4) controls for sample-size differences between models.}

\item The data cleaning and preprocessing steps are clearly and fully described, either in text or as a code pipeline.\\
$\yes$ \textbf{Yes} \quad $\no$ No\\
{\small Methods \S Input data and \S Tasks. No preprocessing applied to abstracts --- used verbatim. Supplementary Section S4 documents exact API payloads for all 7 endpoints.}

\item Instances of combining data from multiple sources are clearly identified, and potential issues mitigated.\\
$\no$ Yes \quad $\yes$ \textbf{No}\\
{\small Not applicable --- all data generated from a single experimental protocol. No combining of external datasets.}
\end{enumerate}

%%=============================================================
\section*{3. Model and training}

\begin{enumerate}[label=\Alph*.,leftmargin=2em]
\item What model architecture is the current model based on?\\
{\small N/A --- We do not propose a new model. We evaluate 8 existing LLMs at inference time: LLaMA 3 8B, Mistral 7B, Gemma 2 9B (Transformer decoders, open-weight), GPT-4, Claude Sonnet 4.5, Gemini 2.5 Pro, DeepSeek Chat, Perplexity Sonar (proprietary). Model identities and versions documented in Methods \S Models and infrastructure.}

\item A Model Card is provided.\\
$\no$ Yes \quad $\yes$ \textbf{No}\\
{\small We do not train or release a model. Instead, we provide a ``Run Card'' (provenance record) for each of the 4,104 inference runs, extending the Model Card concept to the inference layer (Methods \S Protocol design).}

\item The model clearly splits data into different sets for training (model selection), validation (hyperparameter optimization), and testing (final evaluation).\\
$\no$ Yes \quad $\yes$ \textbf{No}\\
{\small Not applicable --- inference reproducibility study. No model training was performed.}

\item The method of data splitting is clearly stated.\\
$\no$ Yes \quad $\yes$ \textbf{No}\\
{\small Not applicable --- no data splitting. All 30 abstracts used as inputs for all models under identical conditions.}

\item The data splitting mimics real-world applications.\\
$\no$ Yes \quad $\yes$ \textbf{No}\\
{\small Not applicable --- no data splitting performed.}

\item The data splitting procedure has been chosen to avoid data leakage.\\
$\no$ Yes \quad $\yes$ \textbf{No}\\
{\small Not applicable --- no data splitting performed.}

\item The interpretability of the model has been studied and clearly validated.\\
$\no$ Yes \quad $\yes$ \textbf{No}\\
{\small Not applicable --- we do not develop or train a model. Interpretability of the reproducibility gap is addressed through the quasi-isolation experiment (Results \S Cloud deployment) and sources-of-non-determinism analysis (Methods \S Sources).}
\end{enumerate}

%%=============================================================
\section*{4. Evaluation}

\begin{enumerate}[label=\Alph*.,leftmargin=2em]
\item The performance metrics used are described and justified in the paper.\\
$\yes$ \textbf{Yes} \quad $\no$ No\\
{\small Methods \S Metrics: Exact Match Rate (EMR), Normalized Edit Distance (NED), ROUGE-L F1, BERTScore F1. Three-level reproducibility framework (bitwise $\to$ structural $\to$ semantic) justified in Results \S Semantic analysis.}

\item Cross-validation of the results is included.\\
$\no$ Yes \quad $\yes$ \textbf{No}\\
{\small Not applicable --- no predictive model. Robustness verified via 10,000-resample bootstrap CIs and balanced 10-abstract subsample analysis (Extended Data Table 4).}

\item Community-accepted benchmark datasets/tasks are used for comparisons.\\
$\no$ Yes \quad $\yes$ \textbf{No}\\
{\small Not applicable in the traditional sense. We designed a novel experimental protocol for inference reproducibility assessment. BERTScore and ROUGE-L are community-standard NLP evaluation metrics.}

\item Baseline comparisons to simple/trivial models are provided.\\
$\no$ Yes \quad $\yes$ \textbf{No}\\
{\small Not applicable --- no predictive model. The ``baseline'' is local deployment (near-perfect reproducibility), against which API deployments are compared. The Together AI quasi-isolation (same weights, different infrastructure) serves as a controlled baseline.}

\item Benchmarks with current state-of-the-art are provided.\\
$\no$ Yes \quad $\yes$ \textbf{No}\\
{\small Not applicable --- no predictive task. We compare against prior work on LLM non-determinism (Atil et al.\ 2024, Yuan et al.\ 2025) in the Discussion. Protocol compared feature-by-feature against 5 existing tools (Supplementary Table S2).}

\item Ablation experiments are included.\\
$\yes$ \textbf{Yes} \quad $\no$ No\\
{\small Supplementary Section S6: Protocol Minimality ablation --- systematically removes each field group from the Run Card schema and assesses which audit questions become unanswerable (Tables S3--S4). Temperature sweep (Extended Data Table 3). Chat-format control (Supplementary Section S7).}

\item The model has been tested on a fully independent dataset.\\
$\no$ Yes \quad $\yes$ \textbf{No}\\
{\small Not applicable --- no predictive model. The study tests inference reproducibility across 8 models on the same input corpus.}
\end{enumerate}

%%=============================================================
\section*{5. Computational resources}

\begin{enumerate}[label=\Alph*.,leftmargin=2em]
\item The paper contains information on hardware/computing resources that were used.\\
$\yes$ \textbf{Yes} \quad $\no$ No\\
{\small Methods \S Models and infrastructure: Apple M4 (24 GB RAM) for local models via Ollama v0.15.5. API models accessed via respective cloud endpoints (OpenAI, Anthropic, Google, DeepSeek, Perplexity, Together AI).}

\item The paper includes information on the computational costs in terms of computation time, parallelization or carbon footprint estimates.\\
$\yes$ \textbf{Yes} \quad $\no$ No\\
{\small Methods \S Protocol overhead: $<$1\% logging overhead ($\sim$25 ms per run, $\sim$4 KB storage per run). Extended Data Table 6 provides detailed field-group-level overhead breakdown. Total: 4,104 runs.}
\end{enumerate}

\end{document}
